\documentclass[10pt]{beamer}
\usetheme{metropolis}

\usepackage{amsmath,amssymb}
\usepackage{mathtools}
\usepackage{xeCJK}
\IfFontExistsTF{Noto Serif SC}{\setCJKmainfont{Noto Serif SC}}{\setCJKmainfont{Noto Serif CJK SC}}
\usepackage{tikz-cd}

% 1. 开启标题背景色填充
\metroset{block=fill}

% theorem environments
\newtheorem{proposition}{Proposition}
\newtheorem{exercise}{Exercise}

% notation
\newcommand{\R}{\mathbb{R}}
\newcommand{\N}{\mathbb{N}}
\newcommand{\F}{\mathcal{F}}
\newcommand{\supp}{\operatorname{supp}}
\newcommand{\vX}{\mathfrak{X}}
\newcommand{\D}{\mathcal{D}}

\title{W9\quad Tensor Products}
%\subtitle{Tensor Products}
\date{November 11, 2026}
\author{LXJ}

\begin{document}

\maketitle

%====================================================================
\section{Chapter 2: Tensors and Differential Forms}

\begin{frame}{The tensor product}
	Let $U, V, W$ be finite dimensional real vector spaces, $V^*$ the
	dual space of $V$.
	\medskip

	\begin{definition}[Tensor product]
		Let $F(V, W)$ be the free vector space over $\R$ generated by
		$V \times W$, and let $R(V, W)$ be the subspace generated by
		the elements
		\[
			\begin{aligned}
				(v_1 + v_2, w) - (v_1, w) - (v_2, w) \\
				(v, w_1 + w_2) - (v, w_1) - (v, w_2) \\
				(av, w) - a(v, w), \qquad (v, aw) - a(v, w)
			\end{aligned}
			\quad
			\Bigl( \substack{a \in \R \\[1pt] v, v_1, v_2 \in V \\
				w, w_1, w_2 \in W} \Bigr).
		\]
		The quotient space $F(V, W) / R(V, W)$ is called the
		\emph{tensor product} of $V$ and $W$, denoted $V \otimes W$.
		The coset of $(v, w)$ is denoted $v \otimes w$.
	\end{definition}
\end{frame}

\begin{frame}{Identities in $V \otimes W$}
	Thus we have the following identities in $V \otimes W$:
	\[
		(v_1 + v_2) \otimes w = v_1 \otimes w + v_2 \otimes w,
	\]
	\[
		v \otimes (w_1 + w_2) = v \otimes w_1 + v \otimes w_2,
	\]
	\[
		(av) \otimes w = a(v \otimes w),
		\qquad
		v \otimes (aw) = a(v \otimes w).
	\]
	\medskip
	\begin{block}{Note}
		The tensor product is \emph{bilinear} by construction:
		$\otimes : V \times W \to V \otimes W$ is linear in each
		argument separately.
	\end{block}
\end{frame}

\begin{frame}[fragile]{Property 1: the universal mapping property}
	Let $\varphi$ be the bilinear map
	\[
		\varphi : V \times W \longrightarrow V \otimes W,
		\qquad (v, w) \longmapsto v \otimes w .
	\]
	Then for any vector space $U$ and any bilinear map
	$\ell : V \times W \to U$, there exists a \emph{unique} linear map
	$\widetilde{\ell} : V \otimes W \to U$ s.t.\ the following diagram
	commutes:
	\begin{center}
		\begin{tikzcd}
			V \times W \arrow[r, "\ell"] \arrow[d, "\varphi"']
			& U \\
			V \otimes W \arrow[ru, "\widetilde{\ell}"', dashed]
		\end{tikzcd}
	\end{center}
	The pair $(V \otimes W, \varphi)$ is said to solve the
	\emph{universal mapping problem} for bilinear maps $\ell$ with
	domain $V \times W$.
\end{frame}

\begin{frame}{Uniqueness and further properties}
	Moreover, $(V \otimes W, \varphi)$ are unique in the sense that if
	$(X, \widetilde{\varphi})$,
	$\widetilde{\varphi} : V \times W \to X$ a bilinear map, also
	solves the universal mapping problem, then there is an isomorphism
	$\alpha : V \otimes W \to X$ s.t.\
	$\alpha \circ \varphi = \widetilde{\varphi}$.
	\medskip
	\begin{enumerate}
		\item[(b)] $V \otimes W$ is canonically isomorphic to
		      $W \otimes V$ (both solve the universal mapping
		      property);
		\item[(c)] $V \otimes (W \otimes U)$ is canonically isomorphic
		      to $(V \otimes W) \otimes U$.
	\end{enumerate}
\end{frame}

\end{document}
